\section{Upper speeds at real panels}\label{sec:panels}

This section treats a non-isolated source position and its matched target
position when both places are archimedean.  The finite-place argument is in
Section~\ref{sec:finite}.  The differentiation argument of
Section~\ref{sec:scale} concerns isolated axes and is not used to supply
the estimate at a real panel.

\subsection{The estimate and the order of limits}

Fix one measured coordinate of the exact graph relation
\eqref{eq:input-6}; this coordinate may also be one of the two measured
coordinates of a two-leg link.  Begin with a mixed many-panel
lexicographic face path that
selects every isolated gap and omits one gap in each higher-rank
component.  Let $S$ and $T$ be the noncompact simple rank-one Levi groups
of an omitted source position and its matched target position.  Compact
and ineffective factors are discarded.  Source actions below may be
lifted to the simply connected covering group.  Write $e_i$ for the
coordinate vector dual to the omitted simple root.  Thus translation by
$t e_i$ has simple gap $t$, rather than root-coroot parameter $t$.

\begin{proposition}[The real-panel upper estimate]\label{prop:panels-upper}
There are positive constants $d_R$, determined by each real rank-one
panel type $R$ and unchanged at the opposite type, with the following
property.  Complete the mixed path by translations in the same split
chamber along its missing coordinates, taking their times to infinity
successively and more slowly than the original path.  Additional slower
translations along already selected isolated coordinates are allowed.
If $t$ is the completing simple gap at the selected panel, and $\ell(g)$
is the residual target Cartan gap, then
\begin{equation}\label{eq:input-11}
 d_T\ell(g)\leq d_S t+O_{\mathrm{tight}}(1).
\end{equation}
Here a tight upper error means that the event on which
$d_T\ell(g)-d_St$ tends to positive infinity has zero mass in the
iterated limit, on every fixed input.  One may add offsets fixed during
all preceding path limits and then take further limits in those offsets;
$t$ in \eqref{eq:input-11} is the resulting shifted gap.  Gap signs are
the eventual orientations on the chosen filters.
\end{proposition}

We record precisely the angular information available during this
completion.  At the first, mixed limit the residual missing gap is
finite, and its Cartan gap differs by a bounded quantity from the
corresponding inverse Fourier-index gap.  A completing translation fixed
at this first limit is an exact input precomposition of the link.  It
therefore retains all the previously sharp angular marginals.  At each
subsequent limit the growth and upper vertex-count argument of
Section~\ref{sec:transfer} applies to the enlarged path.  If a previously
selected gap is repeated, maximal vertex count on that face has already
been attained, so the repetition introduces no additional gap.  In the
stated order of limits, the total split parameters still contract the
relevant face radicals.  Thus every resulting link is mixed and
isometric, has its full sharp marginals, and satisfies
\eqref{eq:input-6} with the indicated input reindexing.  The same reasoning
applies to outer translations in the common split torus after a full
path, including the offsets in Proposition~\ref{prop:panels-upper}.
These assertions use the exact precomposition and angular statements of
Sections~\ref{sec:links} and~\ref{sec:transfer}; they require neither the
height calibration nor the classicality conclusion proved later.

The structural classification in Section~\ref{sec:setup} leaves the
orthogonal, unitary, and symplectic groups in the following table as the
real rank-one panel groups.  The exceptional rank-one group
$F_{4(-20)}$ cannot occur as the noncompact simple factor of a panel Levi
in a higher-rank component.  Its complexification would be a component
of a proper absolute Dynkin subdiagram, whereas no proper subdiagram of
an irreducible Dynkin diagram has a component of type $F_4$.  Equivalently,
in the real restricted-root list a multipliable simple root arising in
such a Levi has double-root multiplicity $1$ or $3$, and not $7$.
The reduced low-rank isomorphisms are accounted for in the table's
orthogonal line.

\subsection{Optimal modules for a real rank-one panel}

Let $R$ be a noncompact rank-one panel factor, and let $H_R$ be its
split generator, normalized so that the simple restricted gap is one.
For a finite-dimensional complex representation $V$ of its real Lie
algebra, extended to the complexification, write
\[
 V_t=\ker(d\pi(H_R)-t),\qquad
 h(V)=\sum_{t>0}\dim V_t,\qquad z(V)=\dim V_0.
\]
For a nontrivial irreducible $V$, the ratio $\dim V/h(V)$ is the
inverse proportion of positive-weight vectors.  We will extract finite-dimensional
source representations from the measured graph relation and compare their
stage dimensions with those of a target module.  A sharp ratio and a uniform
gap away from its maximizers will bound the dimension lost when only
maximizing source summands are retained.  Their largest split weights then
bound source growth.  These are the different roles of the constants
$c_R$ and $d_R$ below.

We allow representations of the simply connected covering group. Thus
the weight lattices below are the full weight lattices, and spin
representations are included. Isomorphic rank-one groups are named only
once: $SU(1,1)$ uses the $SO(2,1)$ line and $Sp(1,1)$ uses the
$SO(4,1)$ line, with the reduced restricted root in both cases.

\begin{lemma}[Optimal panel modules]\label{lem:panels-optimal-modules}
Put $k=n+1$, and let $E$ denote the standard complex module for the
group in the relevant line. Among the nontrivial irreducible complex
Lie-algebra modules, the maximum
\[
 c_R=\max_V\frac{\dim V}{h(V)}
\]
and all modules attaining it are as follows:
\[
\begin{array}{ccl}
 R&c_R&\text{maximizing modules}\\ \hline
 SO(n,1),\ n\geq2&k&E\\
 SU(n,1),\ n>2&k&E,E^*\\
 SU(2,1)&3&E,E^*,\operatorname{Sym}^2E,
                         \operatorname{Sym}^2E^*\\
 Sp(n,1),\ n>2&k&E\\
 Sp(2,1)&7/2&\bigwedge^3_0E.
\end{array}
\]
In the $Sp(2,1)$ row, $\bigwedge^3_0E$ is the primitive exterior cube,
the kernel of symplectic contraction $\bigwedge^3E\to E$.
There is a number $\varepsilon_R>0$ such that every nontrivial
irreducible outside the displayed list satisfies
\[
 \frac{\dim V}{h(V)}\leq c_R-\varepsilon_R.
\]
Every displayed module has integral split weights and a nonzero
split-weight-zero space. It has a preserved nondegenerate Hermitian
form $q$, positive definite on $V_0$, whose matrix in a compact-form
Hilbert norm satisfies
\[
 Q=Q^*=Q^{-1},\qquad d\pi(H_R)^*=d\pi(H_R).
\]
The form pairs $V_t$ perfectly with $V_{-t}$ and pairs it trivially with
all other weight spaces. In particular, for
$V_{>0}=\bigoplus_{t>0}V_t$ and
$V_{\geq0}=\bigoplus_{t\geq0}V_t$,
\[
 V_{>0}^{\perp_q}=V_{\geq0},\qquad
 V_{\geq0}^{\perp_q}=V_{>0},
\]
and the induced form on $V_{\geq0}/V_{>0}$ is positive definite.
These two spaces are algebraic parabolic stages. The largest split
weight occurring among the maximizing modules is
\[
 d_R=\begin{cases}
 2,&R=SU(2,1)\text{ or }Sp(2,1),\\
 1,&\text{otherwise}.
 \end{cases}
\]
The same constants apply to the opposite orientation of the panel.
\end{lemma}

We use this complete representation-theoretic package first; its proof is
given in Subsection~\ref{subsec:panels-module-proofs}.  The next step converts the
residue graph relation into compact operators preserving the two parabolic
stages.  Their finite-dimensional spectral ranges will make the ratio
$c_R$ available for a trace comparison.

\subsection{Algebraic stages and compact graph operators}

Choose a target maximizing module $V$ whose largest split weight is
$d_T$, and put $n_0=\dim V$.  If a compact diagram symmetry interchanges
inequivalent modules, replace $V$ by the direct sum of its diagram
conjugates.  Every summand has the same optimal ratio, so still
\begin{equation}\label{eq:panels-target-ratio}
 h(V)=n_0/c_T.
\end{equation}
Choose the compact-form Hilbert norms and preserved Hermitian forms
provided by Lemma~\ref{lem:panels-optimal-modules}, and denote their common
form symmetry by $Q$, with $Q=Q^*=Q^{-1}$.  In compact projectivity frames
there are local, or Borel, lifts $\pi_T(g)$ of target projectivities such
that
\begin{equation}\label{eq:panels-lifts}
 \pi_T(g)^*Q\pi_T(g)=Q,
 \qquad \|\pi_T(g)\|=\exp(d_T\ell(g)).
\end{equation}
Indeed the connected action lifts to the simply connected module.
Compact Cartan factors act unitarily, and compact diagram factors
permute the chosen summands.  The forms can be chosen consistently under
these permutations: on each irreducible their invariant forms are
unique up to a real scalar, and positivity on the nonempty zero-weight
space fixes its sign.  For the coefficient calculations below, fix Borel
lifts $\pi_T(g)$.  Changing a covering lift acts by a scalar of modulus one
on each irreducible summand, preserving the weight planes, the form and norm
identities above, and the limiting top split-weight projector.  The
calculations use these fixed lifts and do not require a multiplication law
for them.

For a Hilbert space $H_0$, consider matrices in
$\operatorname{End}(V)\otimes H_0$.  We use the total and column
Hermitian forms, whose symmetries relative to the Hilbert norm are,
respectively,
\[
 X\longmapsto QXQ,\qquad X\longmapsto QX.
\]
Their adjoints are denoted by $\sharp$ and
$\sharp_{\mathrm{col}}$.  On the start or end of the link, let
$\mathcal L(z)\subset\mathcal D(z)$ be the closed matrix subspaces
whose columns are spectrally in the strictly positive and nonnegative
split-weight planes of the corresponding $P_z$ or $Q_z$ in
\eqref{eq:input-6}.  The label $z$ specifies the source panel endpoint.
These are split-weight stages, not stages defined by positive eigenvalues
of the preserved form.  Parabolic invariance makes the stages depend only
on the directed endpoint, independently of the opposite endpoint used to
choose a split axis.

\begin{lemma}[Stage geometry and polynomial tests]\label{lem:panels-stages}
The stages are mutual annihilators for either matrix form.  Their column
middle quotient $\mathcal D(z)/\mathcal L(z)$ is a positive Hilbert
space, with its quotient topology in compact frames.  On almost every
sharp endpoint fibre, the closed span of $\mathcal L(z)$ over all labels
is the whole matrix space.  All the stage and projector relations used
below are legitimate consequences of the polynomial fullness statement
in Section~\ref{sec:transfer}.
\end{lemma}

\begin{proof}
Opposite nonzero weights pair perfectly for the preserved form, and its
zero-weight space is positive.  Consequently the positive and
nonnegative weight spaces are mutual annihilators and their middle
quotient is positive.  The assertion passes to spectral stages and
Hilbert multiplicities.  Right multiplication by $Q$ preserves a
column stage, so inserting that additional symmetry gives the same
annihilator relation for the total form.

Here is the algebraic content of the required polynomial tests.  In an
individual compact face frame, restrict the surviving highest-line
embedding of the ambient vertex to the rank-one Levi orbit.  The complex
stabilizer of this line is the parabolic consisting of the split
centralizer and positive root groups.  It embeds the corresponding
complex parabolic variety projectively.  The positive and nonnegative
planes of an algebraic covering module are stabilized by that parabolic.
So is its top split-weight plane: the split centralizer preserves it, and
positive-root operators cannot raise the split weight beyond its maximum.
The Pl\"ucker coordinates of all these planes are therefore rational,
homogeneous coordinates on patches of this embedding.  An orthogonal projector is
rational in a spanning matrix and its conjugate, by Gram inversion.
Thus each projector relation is a rational expression in the real
coordinates of the highest line.  Clear its denominators on a patch and
multiply once more by a denominator that vanishes off that patch.  The
resulting polynomial vanishes there as well, so it is a global polynomial
relation.  Multiplication by norm powers makes it homogeneous of an even
degree when necessary.  This additional clearing factor is essential:
no assertion about a rational expression on its zero-denominator set is
being made.

Apply the fibrewise line-span statement of
Section~\ref{sec:transfer} to these homogeneous polynomial relations
with linear Hilbert-vector coefficients.  In particular, a vector
orthogonal to every measured positive stage would be orthogonal to the
positive plane at every ordinary flag.  The latter planes span $V$: their
span is a nonzero invariant subspace on each chosen irreducible summand.
Tensoring with the coefficient Hilbert space and allowing every matrix
column proves dense spanning.  Countably many determining labels and
patch equations suffice first; strong continuity of the label action
then gives every label.
\end{proof}

Disintegrate over the two sharp endpoint faces.  The exact graph
relation, its determining equations, and the strongly continuous source
Levi action hold on almost every pair of fibres, as proved in
Section~\ref{sec:transfer}.  Use the genuine source amplification of
Section~\ref{sec:links} if a projective composition defect is present.
We temporarily work on such a fibre pair, writing its end and start
Hilbert spaces as $I$ and $H'$, with isometry
\[
 J:\overline I\otimes H'\longrightarrow\mathcal K.
\]
We identify $\overline I\otimes H'$ with
$\operatorname{HS}(I,H')$.

\begin{lemma}[A compact graph operator]\label{lem:panels-compact-graph}
For every $\varepsilon>0$ there is a Hilbert--Schmidt operator
\[
 D:\operatorname{End}(V)\otimes I
       \longrightarrow\operatorname{End}(V)\otimes H'
\]
that carries both end stages to the corresponding start stages for every
label, whose total adjoint carries the start stages back, and that reduces
the scalar/traceless matrix splitting.  If $A:I\to H'$ is its scalar
block, then
\begin{equation}\label{eq:panels-compact-approximation}
 a=\|A\|_2^2>0,
 \qquad
 \|DD^\sharp-1\otimes AA^*\|_1\leq\varepsilon a.
\end{equation}
\end{lemma}

\begin{proof}
Set $C(g)=\Ad\pi_T(g)$.  On a relatively compact projectivity window
$U$, take $v\in E(U)\mathcal K$ and define, with $a,b$ indexing matrix
entries rather than rows of $V$,
\[
 D_{ab}=J^*E(C_{ab})v\in\operatorname{HS}(I,H'),
 \qquad A=J^*v.
\]
The entry convention is explicitly
\[
 \langle\xi',D_{ab}\eta\rangle
 =\langle\overline\eta\otimes\xi',J^*E(C_{ab})v\rangle.
\]
For end and start spectral stage projectors $F(u_-)$ and $L(u_+)$,
pairing a coefficient of $(1-L)DF$ tests the link after multiplication
by conjugated entries of
\[
 (1-L(u_+))C(g)F(u_-).
\]
These entries vanish by \eqref{eq:input-6}.  Determine the equations
first on countable labels and compact windows and then apply bounded
output calculus.  This proves the two stage inclusions; total adjunction
and mutual annihilation prove the reverse inclusions.  In particular the
argument uses complex matrix entries and does not require a real
structure on $V$.

Conjugation fixes scalar matrices and preserves traceless matrices.
These two spaces are also orthogonal for the total form, since
$\Tr(QXQ)=\Tr X$.  Thus the splitting reduces $D$ and $D^\sharp$,
and its scalar block is $A$.

To obtain the approximation, first fix one compact window of nonzero
compression.  On it the norms of $C(g)$ and $C(g)^{-1}$ have a common
finite bound.  Subdivide this fixed window into finitely many Borel
windows of arbitrarily small coefficient oscillation.  At least one
subwindow $U$ still has nonzero compression.  Choose a value $C_0$ on it;
it is a total-form isometry with the already fixed norm bound.  Put
$b_U=\|J^*E(U)\|>0$ and choose a unit $v\in
E(U)\mathcal K$ nearly attaining this norm, with
$\sqrt a=\|J^*v\|\geq b_U/2$.  For a coefficient error $f$ supported
on $U$,
\[
 \|J^*E(f)v\|_2\leq b_U\|f\|_\infty
                    \leq2\sqrt a\,\|f\|_\infty.
\]
There are a fixed finite number of matrix coefficients.  Therefore, if
the window oscillation is $\delta$,
$D=C_0\otimes A+F$ with $\|F\|_2\leq C\delta\sqrt a$.
The form symmetries have norm one, and $C_0C_0^\sharp=1$, so the ideal
inequality $\|XY\|_1\leq\|X\|_2\|Y\|_2$ bounds the difference
in \eqref{eq:panels-compact-approximation} by
$C'(\delta+\delta^2)a$, where $C'$ uses only the original compact
window and the matrix size.  Choose $\delta$ accordingly before the
subdivision to obtain the asserted $\varepsilon a$ bound.
\end{proof}

\subsection{Finite-dimensional ranges and optimal source copies}

Let $W$ be a finite-dimensional nonzero spectral range of $DD^\sharp$,
and let $r$ be the dimension of its intersection with a positive endpoint
stage.  Our goal is to prove
\[
 \dim W\leq c_Sr
\]
and to control the dimension outside the optimal source-module copies by
the defect $c_Sr-\dim W$.  We cannot apply the source representation count
directly to $W$: it need not carry source transports preserving the column
form.  The first Riesz cut gives a range nondegenerate for the total form.
A second cut preserves its entire positive middle quotient and makes the
column form nondegenerate.  This is the range on which we construct the
finite-dimensional source representation.

\begin{lemma}[The two Riesz cuts]\label{lem:panels-riesz}
Let $W$ be a nonzero-eigenvalue Riesz range of $DD^\sharp$, grouping
conjugate eigenvalues together, and let $F_0$ be its total-selfadjoint
Riesz projection.  The integer
$r=\dim(W\cap\mathcal L(z))$ is independent of $z$.  There is a
column-selfadjoint finite-rank projection $e_1$ with range $W_1\subset
W$ such that, for $V_z=e_1\mathcal L(z)$,
\begin{equation}\label{eq:panels-riesz-dimension}
 \dim W-\dim W_1=2(r-\dim V_z).
\end{equation}
The $V_z$ span $W_1$, are isotropic, and have positive perpendicular
quotient.  A finite-dimensional representation of the simply connected
source group on $W_1$ preserves the column form and carries $V_z$ onto
$V_{sz}$.
\end{lemma}

\begin{proof}
Every resolvent away from the nonzero spectrum of a compact operator
preserving a closed stage also preserves that stage.  In fact the
restriction of $\lambda-K$ to the stage is a scalar multiple of identity
minus a compact operator, has Fredholm index zero, and is injective when
the ambient operator is invertible.  It is therefore onto the stage.
Contour integration shows that $F_0$ preserves both stages.  Because
$DD^\sharp$ is total-selfadjoint, grouping conjugate eigenvalues makes
$F_0$ total-selfadjoint.  Its range is consequently nondegenerate, with
the induced stages still mutual annihilators.

Let $P_z^{\mathrm{stage}}$ denote the Hilbert orthogonal projector onto
the matrix stage $\mathcal L(z)$.  The operator $F_0P_z^{\mathrm{stage}}$
is an idempotent with image $W\cap\mathcal L(z)$.  It is norm-continuous
in $z$: $P_z^{\mathrm{stage}}$ and its adjoint are strongly continuous,
and $F_0$ has finite rank.  Hence its rank is locally constant and therefore
constant on the connected rank-one boundary.  The same reasoning will
apply to the second cut below.

Annihilation implies that $F_0^{\sharp_{\mathrm{col}}}$ preserves both
stages as well.  On the positive Hilbert quotient
$\mathcal D(z)/\mathcal L(z)$, let $f$ be the induced idempotent of
$F_0$.  Its image has dimension
\[
 \dim(W\cap\mathcal D(z))-
 \dim(W\cap\mathcal L(z))=\dim W-2r.
\]
Column adjunction descends to Hilbert adjunction on this quotient, so
$F_0F_0^{\sharp_{\mathrm{col}}}$ induces $ff^*$.  Choose $e_1$ to be
the Riesz projection of $F_0F_0^{\sharp_{\mathrm{col}}}$ onto all of
its strictly positive real eigenvalues.  Only finite-rank spectra away
from zero occur.  The projection is column-selfadjoint, preserves the
stages, and has range in $W$.  Passage to the quotient retains exactly
the positive spectral subspace of $ff^*$, whose dimension is $\rank f$.
On $W_1$ the corresponding dimension is
$\dim W_1-2\dim(e_1\mathcal L(z))$.  Thus the second cut retains the
whole original middle quotient: the nonzero spectral range of $ff^*$ is
$\ran f$.  The dimensions lost in $W$ are twice those lost from its
isotropic stage, proving \eqref{eq:panels-riesz-dimension}.

Since $e_1$ is column-selfadjoint, $W_1$ is column-nondegenerate;
$V_z=W_1\cap\mathcal L(z)$ and its perpendicular in $W_1$ is
$W_1\cap\mathcal D(z)$.  Positivity of the middle quotient is inherited
from the ambient one.  Applying $e_1$ to the dense stage span proves that
the $V_z$ span $W_1$.

It remains to produce actual finite-dimensional transports.  Denote the
genuine, strongly continuous ambient source transport by $U_s$; it is a
column-form isometry.  For $s$ near the identity,
\[
 A_s=e_1U_s|_{W_1}
\]
is invertible, by finite rank and strong continuity.  It maps $V_z$ into
$V_{sz}$ for every $z$, and equal constant dimensions make this equality
for every $z$.  Its column adjoint is
$e_1U_{s^{-1}}|_{W_1}$ and carries $V_{sz}$ back onto $V_z$.
Consequently $B_s=A_s^{\sharp_{\mathrm{col}}}A_s$ preserves every
$V_z$ and is near identity.  Its analytic inverse square root is
column-selfadjoint and also preserves every $V_z$.  Thus
\[
 A_s B_s^{-1/2}
\]
is an exact column-form isometry transporting all these planes
simultaneously.  A multiplication rule for the compressions is not
needed.

The pairs $(s,A)$ such that $A$ is a column-form isometry of $W_1$ and
$AV_z=V_{sz}$ for all $z$ form a closed subgroup of
$S\times U(q|_{W_1})$.  Closedness follows from continuity of the plane
family.  The preceding construction supplies local lifts; hence this
Lie group projects onto connected $S$ and its Lie algebra projects onto
$\mathfrak s$.  Its solvable radical maps to zero, since $\mathfrak s$
is simple.  A Levi subalgebra maps onto $\mathfrak s$, and a complementary
simple ideal in that Levi subalgebra maps isomorphically onto
$\mathfrak s$; these are the usual Levi decomposition and complete
reducibility facts for finite-dimensional Lie algebras
\cite[Appendix~B, Lemma~B.1 and Theorem~B.2]{Knapp2023}. Integrating this Lie-algebra section on the simply
connected source group yields the asserted representation.
\end{proof}

\begin{lemma}[The defect and the good subspace]\label{lem:panels-defect}
For a range $W$ of Lemma~\ref{lem:panels-riesz}, put
\[
 \delta_0=c_Sr-\dim W.
\]
Then $\delta_0\geq0$.  There is a sum $G_W\subset W_1$ of optimal
source-module copies whose codimension \emph{in $W$} is at most
$C\delta_0$.  Each copy takes its positive and nonnegative source
weight planes into the measured positive and nonnegative matrix stages
at every label.  The constant depends only on the finite list of panel
types.
\end{lemma}

\begin{proof}
Use the finite-dimensional source representation from
Lemma~\ref{lem:panels-riesz}.  Fix a split endpoint $z_0$, write $h$ for
the positive split-weight dimension of $W_1$, and let $r_0$ be the
negative index of the form on its zero-weight space.  Nonzero opposite
weight spaces pair perfectly.  Since the isotropic $V_{z_0}$ has positive
perpendicular quotient, it exhausts the negative index, and
\[
 \dim V_{z_0}=h+r_0.
\]
The source parabolic at $z_0$ preserves this plane, so in particular it
has a split-weight decomposition.  Its zero-weight part is isotropic in
the zero-weight space and has dimension at most $r_0$.

Let $u$ be the dimension of the trivial summand of $W_1$.  Projection of
$V_z$ to this summand is independent of $z$, by transport, and is the
whole summand because the planes span $W_1$.  At $z_0$ it is the image of
the zero-weight part, so $u\leq r_0$.  Complete reducibility and
\eqref{eq:panels-riesz-dimension} give, with nontrivial irreducibles
listed with repetitions,
\begin{equation}\label{eq:panels-defect-identity}
 \delta_0=(c_S-2)(r-\dim V_{z_0})+(c_Sr_0-u)
       +\sum_{\sigma\ne1}(c_Sh(V_\sigma)-\dim V_\sigma).
\end{equation}
Every term is nonnegative.  Here $c_S>2$, the first dimension difference
is nonnegative, and $c_Sr_0-u\geq(c_S-1)r_0$.

The uniform ratio gap in Lemma~\ref{lem:panels-optimal-modules} now has a
quantitative use.  For a non-optimal $\sigma$, the ratio bound gives explicitly
\[
 c_Sh(V_\sigma)-\dim V_\sigma
 \geq\frac{\varepsilon_S}{c_S-\varepsilon_S}\dim V_\sigma.
\]
Thus its contribution to \eqref{eq:panels-defect-identity} controls its
whole dimension.  The dimension outside $W_1$, the entire
non-optimal and trivial dimension in $W_1$, and $r_0$ are each at most
$C\delta_0$.

Let $O$ be the optimal isotypic sum and $N$ its invariant complementary
sum.  This is an orthogonal decomposition for the preserved form.  Indeed
each optimal irreducible has a nondegenerate invariant Hermitian form,
so its Hermitian contragredient is itself; invariant pairings with
inequivalent non-optimal or trivial constituents vanish.  Put
$U=\operatorname{pr}_O V_{z_0}$.  It is parabolic-invariant and
$\dim U_0\leq r_0$.

We bound its negative-weight part explicitly.  Choose a nonzero real
simple restricted-root vector and its Cartan reflection.  Their bracket
lies in the split line: it lies in the centralizer and changes sign under
the Cartan involution.  After rescaling they form an
$\mathfrak{sl}_2$-triple with diagonal element $2H$, where $H$ is the
simple-gap-one generator.  Every optimal weight for $H$ is integral.
In a finite-dimensional $\mathfrak{sl}_2$-module, raising from weight
$-j$ to zero is injective for each integer $j>0$.  Parabolic invariance
therefore gives injections $U_{-j}\to U_0$.  The finite optimal list has
at most $d_S$ distinct negative integral weights, whence
\[
 \dim U_{\leq0}\leq(d_S+1)r_0.
\]

Since $\dim V_{z_0}\geq h(O)$ and $\dim N\leq C\delta_0$, the
positive part of $V_{z_0}$ has dimension at least
$h(O)-\dim N-(d_S+1)r_0$.  Intersecting it with $O_{>0}$ loses at most
another $\dim N$.  Thus
\[
 \operatorname{codim}_{O_{>0}}(O_{>0}\cap V_{z_0})
 \leq2\dim N+(d_S+1)r_0\leq C\delta_0.
\]
Moreover pairing $O_{\geq0}$ against $V_{z_0}$ only tests $U_{\leq0}$,
so
\[
 \operatorname{codim}_{O_{\geq0}}
        (O_{\geq0}\cap V_{z_0}^{\perp})\leq C\delta_0.
\]

Write $O=\bigoplus_\sigma V_\sigma\otimes M_\sigma$ over the finite
optimal list.  For each basis vector of $(V_\sigma)_{>0}$ impose on
$m\in M_\sigma$ the condition $v\otimes m\in V_{z_0}$; for each
basis vector of $(V_\sigma)_{\geq0}$ impose
$v\otimes m\in V_{z_0}^{\perp}$.  Each condition has rank at most the
corresponding codimension just bounded.  There are a fixed finite number
of basis vectors, and their module dimensions are fixed.  The resulting
multiplicity subspaces therefore remove at most $C\delta_0$ dimensions
from $O$.  Their full-module sum $G_W$ has the claimed codimension in
$W$.  Transport gives both plane conditions at all labels.  Finally,
$V_z\subset\mathcal L(z)$, and its perpendicular in $W_1$ equals
$W_1\cap\mathcal D(z)$ because $e_1$ is column-selfadjoint.  Thus the
nonnegative condition also lifts to the ambient measured stage.
\end{proof}

The defect estimate now connects the finite-dimensional source ranges to
the target trace.  The compact graph operator is nearly a scalar
multiplicity operator.  Its spectral multiplicities, weighted by eigenvalue
absolute value, have total at least $n_0^2a$ up to a small error, while the
corresponding positive-stage total is at most $n_0^2a/c_T$ up to that error.  Comparing these two quantities first proves equality of the panel
ratios.  The remaining defect is then too small to exclude every scalar
vector from the optimal source copies.

\begin{lemma}[Ratio equality and a scalar witness]\label{lem:panels-scalar-witness}
Matched real panel types satisfy $c_S=c_T$.  On almost every endpoint
fibre there is a nonzero vector $\eta$ such that
$V\otimes\eta$ is contained in the algebraic span of ranges of maps
\[
 B:V_\sigma\longrightarrow V\otimes H'
\]
over optimal source irreducibles $V_\sigma$, where $B$ carries both
source weight stages into the measured column stages at every label.
The analogous assertion holds at the end of the link.
\end{lemma}

\begin{proof}
Sum the Riesz ranges $W$ of $DD^\sharp$ with weight $w$ equal to the
absolute value of the associated eigenvalue, grouping conjugates with
their common absolute value.  Lidskii's theorem \cite{Simon2005} and
\eqref{eq:panels-compact-approximation} give
\[
 \sum_Ww\dim W\geq(n_0^2-\varepsilon)a.
\]
For the restriction to a fixed invariant positive stage, the sum of
absolute eigenvalues, with their Riesz multiplicities $r$, is bounded
above by its trace norm.  The compression of $1\otimes AA^*$ to that
stage is positive.  If $P_z^{\mathrm{stage}}$ is the Hilbert stage projector, its partial
matrix trace is
\[
 \Tr_{\operatorname{End}(V)}P_z^{\mathrm{stage}}
   =n_0h(V)\,1_{H'}.
\]
Indeed each spectral positive-plane projector has rank $h(V)$, and there
are $n_0$ matrix columns.  The compression therefore has trace
$a n_0h(V)=a n_0^2/c_T$.  Compressing the error in
\eqref{eq:panels-compact-approximation} does not increase its trace
norm.  Consequently
\begin{equation}\label{eq:panels-weighted-ranks}
 \sum_Ww r\leq(n_0^2/c_T+\varepsilon)a.
\end{equation}
These restriction multiplicities are exactly
$\dim(W\cap\mathcal L(z))$: the resolvent argument in
Lemma~\ref{lem:panels-riesz} also identifies the restricted Riesz
projections.  The inequality $\dim W\leq c_Sr$ now gives $c_T\leq c_S$
as $\varepsilon\downarrow0$.  The real non-isolated type set is invariant under the type permutation:
Corollary~\ref{cor:finite-arch-separation} preserves the real types, and
Proposition~\ref{prop:transfer-galleries} preserves the Coxeter components.
Applying this inequality round each of its finite cycles gives equality
on every matched pair.

With equality established, \eqref{eq:panels-weighted-ranks} and the lower
bound above imply
\[
 \sum_W w\delta_0\leq(c_S+1)\varepsilon a.
\]
The scalar block of $DD^\sharp$ is exactly $AA^*$.  Put
\[
 S_W=W\cap(\mathbb C1_V\otimes H').
\]
The nonzero such subspaces are precisely the eigenspaces
$1_V\otimes\ker(AA^*-\mu)$, $\mu>0$, so
$\sum_Ww\dim S_W=\Tr(AA^*)=a$.  If $G_W\cap S_W=\{0\}$ for every
$W$, the quotient map $S_W\to W/G_W$ would be injective.  Thus
$\dim S_W\leq\operatorname{codim}_W G_W$, and
Lemma~\ref{lem:panels-defect} would give
\[
 a\leq C\sum_Ww\delta_0\leq C(c_S+1)\varepsilon a,
\]
which is impossible for sufficiently small $\varepsilon$.  Hence one
$G_W$ contains a nonzero scalar matrix $1_V\otimes\eta$.  Projecting
its optimal module copies to each matrix column gives the maps $B$ in
the statement.  The columns of this scalar matrix are
$v_i\otimes\eta$, proving the assertion for every vector of $V$.
Interchanging start and end proves the other assertion.

If a genuine source amplification was used, apply a coefficient
functional in the added Hilbert factor to the maps $B$ and to the scalar
witness.  At least one coefficient of a nonzero witness is nonzero, and
all plane equations are unchanged because the added factor is a
multiplicity for those equations.  This yields a witness on an original
fibre.
\end{proof}

\subsection{Global section domains and their positive matrix action}

The scalar witness is so far only nonzero on almost every endpoint fibre.  We now
pass from those witnesses to a dense global domain, and transfer the source
matrix algebras to it.  The two stage conditions make all cross pairings of
sections source intertwiners.  Positivity on source zero-weight spaces will
then turn the transferred form symmetry into a strictly positive action.
This action will be used on fixed vectors of its algebraic domain in the
speed estimate.

Call a map $B$ in Lemma~\ref{lem:panels-scalar-witness} a \emph{section}.
In a global endpoint space $H_0$, use the same compact frames at the
sharp endpoint, denote the section space for $V_\sigma$ by
$\mathscr S_\sigma$, and put
\[
 \mathcal E=\sum_\sigma\sum_{B\in\mathscr S_\sigma}B(V_\sigma),
 \qquad
 \mathscr D=\{\eta\in H_0:V\otimes\eta\subset\mathcal E\}.
\]
All sums defining $\mathcal E$ are algebraic.  In particular no operator
norm on a completed space is being asserted for the action constructed
below.

\begin{lemma}[Density of the scalar domain]\label{lem:panels-global-domain}
The scalar domain $\mathscr D$ is dense in each global endpoint space.
Its defining section equations are decomposable over the sharp endpoint
face, and the projection onto its closure disintegrates to the fibrewise
closure projections.
It is preserved by the applicable base and complementary right
symmetries, by target translations after the compact-frame correction,
and by source panel transports.  It is also preserved by the split
translations that fix all labels of the panel and escape in all factors.
\end{lemma}

\begin{proof}
Fix countably many dense labels and bases of the finitely many optimal
modules.  Each section condition is a bounded linear equation on
$\Hom(V_\sigma,V\otimes H_0)$: multiply the image of a source-plane
projector by the complement of the measured-plane projector.  Thus
$\mathscr S_\sigma$ is a closed common kernel.  One can realize all
conditions as the kernel of a single bounded map to a countable Hilbert
sum by inserting summable positive weights.  Strong continuity recovers
the conditions at every label.

Choose bases $u_{\sigma,a}$ of $V_\sigma$ and $v_i$ of $V$.  Algebraic
membership $v_i\otimes\eta\in\mathcal E$ can be expressed using only
one unknown section $B_{i,\sigma,a}$ per source basis element:
\[
 v_i\otimes\eta=\sum_{\sigma,a}B_{i,\sigma,a}u_{\sigma,a}.
\]
Indeed finite sums of sections of a fixed type may be combined, including
their scalar coefficients.  The simultaneous equations for all $i$
define another closed linear kernel in the product of $H_0$ and finitely
many section spaces.  The scalar domain is its coordinate image in
$H_0$.  Kernel projections and range-support projections of bounded
decomposable maps disintegrate, so the projection onto
$\overline{\mathscr D}$ disintegrates to the fibrewise closure
projections.  Lemma~\ref{lem:panels-scalar-witness} makes this projection
nonzero; no measurable choice of individual witnesses is necessary.

The defining equations commute with base multipliers and all
complementary right symmetries.  The closure projection consequently
belongs to $\mathcal A_j$.  An actual target translation changes compact
angular frames by a projectivity depending on the sharp face, but not on
the panel label.  The corresponding matrix correction $\pi_T$ is
boundedly invertible for each fixed translation: its Levi projection is
bounded in compact frames, and the same is true for its inverse.  Correct
each transported section by this matrix.  The scalar columns remain in
the span of the corrected sections by combining them with the entries of
the inverse matrix.  These entries are own-face spectral multipliers,
which preserve the section equations.  The inverse translation proves
equality of the transported scalar domains.  No correction is needed
for base or complementary right symmetries.  It follows that the closure
projection is fixed by the measured right action and hence belongs to
$\mathcal N$.

Source panel covariance follows directly as well: if $B$ is a section,
then $U_sB\sigma(s)^{-1}$ is a section with the labels reindexed.  Since
$U_s$ acts only on the Hilbert coefficient, it preserves the scalar
domain.  Split translations that fix every panel label preserve its
defining equations without reindexing.  Choose such translations
escaping in every factor, as in the original mixed face path.  The
closure projection $p\in\mathcal N$ is invariant under this sequence.
Mixing in $\mathcal N$ therefore gives
$\tau(p)=\lim\tau(p\sigma(a)p)=\tau(p)^2$.
Since $p\ne0$, it equals $1$.  This proves density, and also makes
$\mathcal E$ dense in $V\otimes H_0$.
\end{proof}

For the following algebraic argument, let $H_{\mathrm{ax}}$ range over the source
conjugates of the normalized split generator.  On an optimal source module
write $\operatorname{sgn}(d\sigma(H_{\mathrm{ax}}))$ for the operator with values
$1,0,-1$ on its positive, zero, and negative weight spaces.  The generator
$H_{\mathrm{ax}}$ specifies an ordered pair of opposite endpoints.  Its positive and
nonnegative stages belong to the attracting endpoint, and its negative
and nonpositive stages belong to the repelling endpoint.  Thus section
equations at both endpoints can test all three sign spaces; no choice of
split axis is inferred from a single endpoint.

\begin{lemma}[The algebra generated by split signs]\label{lem:panels-sign-algebra}
On an optimal source irreducible, the conjugate split sign operators
generate the full complex matrix algebra.  An operator between two
optimal source irreducibles that intertwines all these signs is either
zero or an ordinary source-module isomorphism, unique up to scalar.
\end{lemma}

\begin{proof}
Let $\mathscr A$ be the unital algebra generated by the signs on an
irreducible module $W$.  The connected source normalizes $\mathscr A$.
Its Jacobson radical has a nonzero common kernel: a finite-dimensional
nilpotent algebra annihilates a nonzero vector.  That common kernel is
source-invariant because the radical is characteristic.  Irreducibility
makes it all of $W$, so the radical is zero.  The connected source cannot
permute the finitely many central blocks of the semisimple algebra, and
their images are invariant subspaces.  There is consequently one block,
and
\[
 W=E\otimes F,\qquad
 \mathscr A=\operatorname{End}(E)\otimes1.
\]
Derivations of a matrix algebra are inner.  For each source Lie-algebra
element choose the unique traceless implementer of its derivation on
$\operatorname{End}(E)$.  These implementers form a Lie representation:
the defect in a commutator is both scalar and traceless, hence zero.
Subtracting this representation from the action on $W$ leaves a
representation in the commutant $1\otimes\operatorname{End}(F)$.
This gives a tensor decomposition of the Lie representation.

Suppose first that the complexified source algebra is simple.  Two
nontrivial irreducibles of a simple complex Lie algebra do not have an
irreducible tensor product.  For completeness, if their highest weights
are $\lambda,\mu\ne0$, the tensor product contains the constituent of
highest weight $\lambda+\mu$.  In the Weyl dimension formula each
positive coroot contributes
\[
 \frac{(a+r)(b+r)}{r(a+b+r)}\geq1,
 \quad
 a=\langle\lambda,\alpha^\vee\rangle,\quad
 b=\langle\mu,\alpha^\vee\rangle,\quad
 r=\langle\rho,\alpha^\vee\rangle>0.
\]
There is a positive coroot pairing positively with both nonzero dominant
weights, because the root system is irreducible.  Its factor is strictly
larger than one.  Hence
$\dim V_\lambda\dim V_\mu>\dim V_{\lambda+\mu}$, excluding
irreducibility.  In the displayed tensor decomposition, irreducibility
of $W$ therefore makes one factor trivial.  The sign is nonscalar, so
the $E$ factor is nontrivial; a trivial $F$ has dimension one by
irreducibility.  This proves $\mathscr A=\operatorname{End}(W)$.

The remaining real-simple case is $\mathfrak{so}(3,1)$, whose
complexification has two simple summands.  Its only optimal module is
the Lorentz vector module.  Its split weights are $1,0,0,-1$, so a split
sign is the split generator itself.  The linear span of its source
conjugates is the image of a nonzero adjoint ideal in the real-simple
algebra and is therefore the whole Lie-algebra image.  The associative
algebra it generates acts irreducibly and equals the full matrix algebra
by the same radical and single-block argument, or by Burnside's theorem.

An intertwiner of all signs has invariant kernel and image for their
full matrix algebras.  A nonzero such map is therefore an isomorphism,
and the space of such maps has dimension one.  Conjugating by the source
preserves this line, because it permutes the axes.  Connected
semisimplicity has no nontrivial one-dimensional characters, so the line
is fixed pointwise.  The map is an ordinary source intertwiner.
\end{proof}

\begin{proposition}[The positive algebraic action]\label{prop:panels-positive-action}
On $\mathcal E$ there is a well-defined action of the finite-dimensional
algebra $\bigoplus_\sigma\operatorname{End}(V_\sigma)$ given by
\[
 R(X)(Bv)=BX_\sigma v.
\]
It is adjoint for the target column form according to source-form
adjunction.  If $\mathbf Q=(Q_\sigma)_\sigma$ is the tuple of source
form symmetries, then
\begin{equation}\label{eq:input-12}
 q(x,R(\mathbf Q)x)>0\qquad(0\ne x\in\mathcal E).
\end{equation}
\end{proposition}

\begin{proof}
For two sections $B:V_\sigma\to V\otimes H_0$ and
$C:V_\tau\to V\otimes H_0$, represent their cross pairing by
$K:V_\tau\to V_\sigma$ using the nondegenerate source form:
\[
 q(Bv,Cw)=q_\sigma(v,Kw).
\]
If $w$ is positive at an endpoint and $v$ is nonnegative, their images
belong to mutually annihilating measured stages.  Hence $Kw$ is positive
at that endpoint.  Interchanging positive and nonnegative gives the
corresponding nonnegative inclusion.  Repeating at the opposite endpoint
gives the negative and nonpositive inclusions, and hence exact
preservation of the three split signs.  Lemma~\ref{lem:panels-sign-algebra}
therefore makes $K$ an ordinary source intertwiner.  Schur's lemma shows
that inequivalent types are orthogonal and, within a fixed type,
\[
 q(Bv,Cw)=q_\sigma(v,w)\,b_\sigma(B,C)
\]
for a Hermitian multiplicity pairing $b_\sigma$.

This pairing is positive semidefinite.  Choose a nonzero source
zero-weight vector $v$; its form is positive by
Lemma~\ref{lem:panels-optimal-modules}.  The vector $Bv$ lies in the
measured nonnegative stage, whose form is nonnegative.  Thus
$b_\sigma(B,B)\geq0$, and the same argument applied to arbitrary
linear combinations proves matrix positivity for every finite family of
sections.

To justify descent rather than assume it, consider the algebraic
evaluation map
\[
 T:\bigoplus_\sigma V_\sigma\otimes\mathscr S_\sigma
       \longrightarrow\mathcal E,
 \qquad T(v\otimes B)=Bv.
\]
Its pullback form is the direct sum of $q_\sigma\otimes b_\sigma$.
For a finite collection of sections the radical is precisely the sum of
$V_\sigma$ tensored with the radicals of the positive semidefinite
multiplicity forms.  A radical vector remains radical after any further
finite collection is added: Cauchy--Schwarz for a positive semidefinite
form makes its pairing with every newly added vector zero.  Its
evaluation therefore pairs to zero with all of $\mathcal E$, and is
zero by the density in Lemma~\ref{lem:panels-global-domain} and ambient
nondegeneracy.  Conversely a vector evaluating to zero is in this
radical, by pairing with the same finite module span.  Thus the kernel
of $T$ is exactly the radical.  It is invariant under every source
matrix action, proving well-definedness and the adjoint assertion.

Finally $q_\sigma(v,Q_\sigma w)$ is the compact-form positive Hilbert
pairing.  The pullback of $q(x,R(\mathbf Q)y)$ is consequently the
tensor product of this positive pairing with $b_\sigma$, summed over
types.  Its only null vectors are the radicals just identified with
$\ker T$.  This proves the strict inequality
\eqref{eq:input-12} for every nonzero evaluated vector.
\end{proof}

\subsection{Section transport and the excess estimate}

For $\eta\in\mathscr D$, let $Z_\eta(X)$ be the $n_0\times n_0$
matrix with Hilbert-vector entries whose $i$th column is
$R(X)(v_i\otimes\eta)$.  In particular $Z_\eta(1)=1_V\otimes\eta$.
For fixed $\eta$, this is a linear map from a fixed finite-dimensional
matrix algebra into a Hilbert space.  Choosing finite section
expressions for its columns gives a finite bound on each matrix unit;
no bounded extension of $R$ beyond its algebraic domain is implied.

Let $J_t$ be a slower translated link as in
Proposition~\ref{prop:panels-upper}.  At the selected panel the exact
graph relation is reindexed across its ends by the source split
translation $s_t$, of simple gap $t$.  Choose its orientation so that
the formulas below use $\sigma(s_t)$; reversing the axis gives the
opposite convention.  Other completing shifts do not change this
panel's indexing.  All frames are taken at the previously sharp face
slots, whose laws were retained by the exact precomposition.

\begin{lemma}[Action exchange and its fixed-input bound]\label{lem:panels-action-exchange}
Construct $D_t(v)$ from $J_t$ on a compact projectivity window as in
Lemma~\ref{lem:panels-compact-graph}, without requiring its approximation
property, and let $A_t(v)=J_t^*v$ be its scalar block.  If $\eta$ belongs
to the end scalar domain, then $A_t(v)\eta$ belongs to the start scalar
domain and
\begin{equation}\label{eq:input-13}
 D_t(v)Z_\eta(X)=
 Z_{A_t(v)\eta}
 \bigl((\Ad\sigma(s_t)X_\sigma)_\sigma\bigr).
\end{equation}
For fixed scalar-domain inputs at the two ends, weighting $v$ by
$e^{-2d_T\ell(g)}$ and cutting to $d_T\ell(g)>d_St+M$ bounds every
paired matrix entry in this identity by
$C e^{-2M}\|v\|$.  The constant is independent of $t$, $M$, and the
location of the compact window.
\end{lemma}

\begin{proof}
Write $D_t^{lm}$ for the block from input matrix column $m$ to output
matrix column $l$.  The stage calculation in
Lemma~\ref{lem:panels-compact-graph} shows that, for an end section $B$,
\[
 D_t^{lm}(v)B\sigma(s_t)^{-1}
\]
is a start section.  This gives action exchange on each column block,
with $X_\sigma$ replaced by
$\Ad\sigma(s_t)X_\sigma$.  Applying the statement with $X=1$ to the
scalar columns and using scalar/traceless reduction yields
\[
 D_t(v)(1_V\otimes\eta)=1_V\otimes A_t(v)\eta.
\]
Every output column is a sum of section evaluations, proving that
$A_t(v)\eta$ is in the start scalar domain.  The general column-block
identity then gives \eqref{eq:input-13} by summation.

To estimate a fixed matrix entry, pair its output against a fixed
start scalar-domain test $\xi$.  A Hilbert pairing with a column $x$
can be written as a column-form pairing with $Qx$.  Move the algebra
action on the right side of \eqref{eq:input-13} to this fixed test by
Proposition~\ref{prop:panels-positive-action}.  Expand that fixed test
in finitely many section expressions and matrix units.  On the finite
optimal list,
\[
 \|\Ad\sigma(s_t)\|\leq e^{2d_St}\qquad(t\geq0),
\]
so the norm of the moved action on the fixed test is at most
$C_\xi e^{2d_St}$.  If
\[
 f_{t,M}(g)=e^{-2d_T\ell(g)}
             1_{\{d_T\ell(g)>d_St+M\}},
\]
then
\[
 \|A_t(E(f_{t,M})v)\|_2
 \leq\|E(f_{t,M})v\|
 \leq e^{-2d_St-2M}\|v\|,
\]
because $J_t$ is an isometry.  The remaining fixed input contributes
only its Hilbert norm.  This proves the bound.  In particular the proof
estimates the action on a fixed test vector and never the norm of an
unbounded action on the moving vector $A_t(v)\eta$.
\end{proof}

We spell out the conjugation convention needed to read this estimate as
an output statement.  Under
$\overline I\otimes H'\simeq\operatorname{HS}(I,H')$, pairing an
entry
\[
 \sum_b (J_t^*E(C_{ab})v)Z_b
\]
against $\xi$ pairs $v$ against
\[
 \sum_b E(\overline{C_{ab}})
                J_t(\overline{Z_b}\otimes\xi).
\]
Conjugating the vector test replaces $J_t$ by its conjugate link
$\overline J_t$.  Thus Lemma~\ref{lem:panels-action-exchange}, by
duality and exhaustion of compact windows, says that the matrix with
entries
\begin{equation}\label{eq:panels-conjugate-bound}
 \bigl[e^{-2d_T\ell(g)}\Ad\pi_T(g)\bigr]\,
       \overline J_t\bigl(Z_\eta(X)\otimes\overline\xi\bigr)
\end{equation}
has norm at most $Ce^{-2M}$ on the excess cutoff.  Here the matrix
coefficients of $\Ad\pi_T(g)$ act by output spectral calculus.  The
normalized operator has norm at most one, since
$\|\pi_T(g)^{-1}\|=\|\pi_T(g)\|$.  Hence compact-window exhaustion
is valid; the constant in the estimate has no window dependence.

\begin{proof}[Proof of Proposition~\ref{prop:panels-upper}]
Take the iterated path limits in their stated order and then the limit
$M\to\infty$.  Equation~\eqref{eq:panels-conjugate-bound} vanishes on
the positive infinite-upper-excess summand.  Put
\[
 b(g)=e^{-d_T\ell(g)}\pi_T(g).
\]
Both $b$ and $b^*$ are contractions.  Right multiplication of the
vanishing matrix by $Q$ uses
\[
 \pi_T(g)^{-1}Q=Q\pi_T(g)^*.
\]
It changes normalized conjugation into left and right multiplication on
$ZQ$ by $b,b^*$.  Further multiplication on the outside by $b^*,b$
gives the vanishing matrix with both outside factors $b^*b$.

Compact Cartan decomposition identifies their limit.  In the surviving
Levi block, if $\pi_T(g)=k_1\exp(\ell(g)H_T)k_2$ in compact frames,
then
\[
 b^*b=k_2^*\exp\bigl(2\ell(g)(H_T-d_T)\bigr)k_2
\]
tends, on divergent residual gap, to the projector on the highest input
split-weight plane.  Denote this projector by $\Pi$; as a function of
the additional reverse angular vertex, it is the $Q$-image of the highest
plane for the inverse projectivity.  Compact factors commute with $Q$,
so these descriptions agree.

All frame passages can be made on continuity sets of arbitrarily large
fixed input mass.  The sharp endpoint base laws are exactly the same
throughout the slower translations, because those translations stabilize
the bases.  Approximate the Borel compact frames on such sets and then
remove their complements.  The additional reverse vertex is sharp by the
total angular statements recalled at the beginning of the section.
Thus $\Pi$ is available as a sharp input spectral projector in the
limit, without a continuity assumption on a moving Borel frame.

Suppose that the infinite-excess summand is nonzero.  Fixed target shifts
at either input alter inverse-index gaps only by bounded amounts, so
they preserve this summand.  The independent all-factor contracted
symmetries preserve it as well, in the same order of limits as for the
mixed link.  The scalar-mass argument of Section~\ref{sec:links}
therefore makes its compression a positive scalar multiple of identity.
Write its scalar mass as $m>0$.  The restricted link divided by
$\sqrt m$ is an isometry and retains the sharp angular marginals.
Indeed the angular subeffect at a sharp input projection $P$ is supported
on $P$, while its complementary cut is supported on $1-P$; their sum is
$mI$, so the first subeffect equals $mP$.  This is the sharp-compression
criterion of Proposition~\ref{prop:links-link}.  We may therefore pull the
additional reverse vertex, denoted $Q_y$, through the restricted conjugate
link.
Pulling the two outside copies of $\Pi$ evaluates both at the same end
slot, and removal of the scaled isometry and a fixed nonzero start test
$\xi$ gives
\begin{equation}\label{eq:panels-projector-vanishing}
 \left[\Pi(v)Z_\eta(X)Q\Pi(v)\right]_{v=Q_y}=0.
\end{equation}
The spectral multipliers in this formula act on its Hilbert-vector
entries.  In particular no section action has been moved through either
projector.

To apply fibrewise polynomial fullness, we need this identity at every
panel label with $\eta$ and $X$ fixed.  Conjugate the completing axis by a
compact source Levi rotation.  This conjugates the entire
path, moving fixed source transports into its two inputs, and preserves
the same scalar excess mass.  Its only effect on the section argument
is a fixed conjugation of the source reindexing.  The scalar domains are
preserved by these transports by
Lemma~\ref{lem:panels-global-domain}.  To see explicitly that the data can be held fixed, the section transport
is $B\mapsto U_k B\sigma(k)^{-1}$, whence
$U_kR(X)U_k^{-1}=R(\Ad\sigma(k)X)$. For prescribed $\eta,X$ apply
the unrotated argument to $\eta'=U_k^{-1}\eta$ and
$X'=(\Ad\sigma(k))^{-1}X$. Then
\[
 (1\otimes U_k)Z_{\eta'}(X')=Z_\eta(X).
\]
Thus Equation~\eqref{eq:panels-projector-vanishing} holds at every
rotated label for the same fixed $\eta$ and $X$.  Impose countably many dense labels first and
extend by strong label continuity.  Lemma~\ref{lem:panels-stages}, with
its extra denominator clearing, now gives the vanishing at every
ordinary residue value $v$ for fixed $\eta$ and $X$.

Fix such an ordinary value and take a constant nonzero
$w\in\ran\Pi(v)$ in the compact frame.  Choose any nonzero
$\eta\in\mathscr D$ at the end and set $X=\mathbf Q$.  Pairing
\eqref{eq:panels-projector-vanishing} with $w\otimes\eta$ gives
\[
 0=\big\langle w\otimes\eta,
          R(\mathbf Q)(Qw\otimes\eta)\big\rangle
   =q\bigl(Qw\otimes\eta,
          R(\mathbf Q)(Qw\otimes\eta)\bigr).
\]
The last quantity is strictly positive by \eqref{eq:input-12}, since
$Qw\otimes\eta\ne0$.  This contradiction proves that the
infinite-upper-excess mass is zero and establishes
\eqref{eq:input-11}.

Every bound used fixed square-integrable scalar-domain data; these data
are dense by Lemma~\ref{lem:panels-global-domain}, so the zero-mass
conclusion extends to all fixed inputs.  To check the offset assertion, absorb a displacement in this panel's
split direction into the total parameter $t$ before applying
$\|\Ad\sigma(s_t)\|\leq e^{2d_St}$.  For each fixed negative offset,
that total parameter is eventually positive; offsets are exhausted only
after the ordinary path limits.  A split translation $a$ in another
completing direction fixes every label of this panel.  It sends a section
$B$ to $U_aB$, so on $\mathcal E$ one has
\[
 R(X)U_a=U_aR(X).
\]
Its unitary action therefore leaves unchanged the norms of the matrix
actions moved to fixed tests in Lemma~\ref{lem:panels-action-exchange}.
This proves the stated uniformity for other completing translations and
the subsequent offset limits.
\end{proof}

The upper-speed proof is complete subject to the optimal-module package.
We now verify that package without using any of the operator constructions
above.

\subsection{Proof of the optimal-module classification}
\label{subsec:panels-module-proofs}

We prove the weight bounds, the uniform separation, and the assertions
about forms separately. The only nonsimple complexified Lie algebra
that occurs is
$\mathfrak{so}(3,1)_{\mathbb C}\simeq
\mathfrak{sl}_2(\mathbb C)\oplus\mathfrak{sl}_2(\mathbb C)$;
it is included explicitly below.

\begin{lemma}[Weight multiplicity comparisons]
\label{lem:panels-weight-multiplicities}
Let $V$ be a nontrivial irreducible module for a complex semisimple
Lie algebra, let $m_\mu$ be its full-Cartan weight multiplicities, and
let $\alpha_1,\ldots,\alpha_r$ be absolute simple roots. Then
\[
 m_0\leq\sum_{j=1}^r m_{\alpha_j}.
\]
For any root $\alpha$ and any weight $\mu$ satisfying
$\mu(\alpha^\vee)>0$,
\[
 m_\mu\leq m_{\mu-\alpha}.
\]
The multiplicities are constant on absolute Weyl orbits.
\end{lemma}

\begin{proof}
The joint raising map
\[
 V[0]\longrightarrow\bigoplus_{j=1}^r V[\alpha_j],
 \qquad v\longmapsto(E_{\alpha_j}v)_{j=1}^r
\]
is injective. Indeed, a vector in its kernel is a highest-weight
vector of weight zero: the simple raising operators generate the
positive nilpotent algebra. The finite-dimensional module generated
by such a vector is trivial. To see the last assertion directly,
restrict to each simple-root $\mathfrak{sl}_2$: a highest vector of
weight zero is also killed by its lowering operator. Thus all simple
raising and lowering operators, and the Cartan algebra, kill the
vector. A nonzero such vector would be an invariant vector in the
nontrivial irreducible $V$, which is impossible.

For the second assertion, restrict $V$ to the root
$\mathfrak{sl}_2$ corresponding to $\alpha$. In every irreducible
$\mathfrak{sl}_2$ summand, lowering is injective on a positive
$\alpha^\vee$-eigenspace: only the lowest weight is killed, and that
weight is nonpositive. The lowering map is therefore injective on
$V[\mu]$ and takes it into $V[\mu-\alpha]$. Finally, Weyl-group
representatives in the simply connected group give isomorphisms
between the weight spaces in a Weyl orbit.
\end{proof}

For a nonzero full-Cartan orbit $\mathcal O$, let
$p(\mathcal O)$ be the number of its weights that are strictly
positive on $H_R$. For a proposed bound $c$, define its surplus by
\[
 s_c(\mathcal O)=c\,p(\mathcal O)-|\mathcal O|.
\]
If $m_{\mathcal O}$ is its multiplicity, then the exact identity
\begin{equation}\label{eq:panels-surplus}
 c\,h(V)-\dim V
 =\sum_{\mathcal O\neq\{0\}}
       m_{\mathcal O}s_c(\mathcal O)-m_0
\end{equation}
separates the full-Cartan zero weight from the other weights. A weight
that is nonzero on the full Cartan but zero on $H_R$ is already
included in its nonzero orbit in this identity.

\begin{lemma}[Orthogonal weight bound]\label{lem:panels-orthogonal-ratio}
For every nontrivial irreducible complex module $V$ of $SO(n,1)$,
with all covering-group modules allowed,
$\dim V\leq(n+1)h(V)$, with equality precisely for the standard
module.
\end{lemma}

\begin{proof}
First let $n+1=2b+1$ with $b\geq2$. The absolute root system is
$B_b$, in coordinates $e_1,\ldots,e_b$, and the split functional
extracts the first coordinate. The Weyl group acts by all signed
permutations. If a nonzero weight has exactly $j$ nonzero
coordinates, symmetry gives
\[
 \frac{p(\mathcal O)}{|\mathcal O|}=\frac{j}{2b}.
\]
Consequently every nonzero orbit has positive surplus for
$c=2b+1$. The short-root orbit has size $2b$, positive count one,
and surplus one. The long-root orbit has size $2b(b-1)$,
positive count $2(b-1)$, and surplus
\[
 (2b+1)2(b-1)-2b(b-1)=2(b^2-1).
\]
Write $m_{\mathrm s}$ and $m_{\mathrm l}$ for their
multiplicities. There are $b-1$ long absolute simple roots and one
short one, so Lemma~\ref{lem:panels-weight-multiplicities} gives
\[
 m_0\leq(b-1)m_{\mathrm l}+m_{\mathrm s}.
\]
Substituting this bound in Equation~\eqref{eq:panels-surplus}
leaves the nonnegative expression
\[
 \bigl(2(b^2-1)-(b-1)\bigr)m_{\mathrm l}
 +\sum_{\substack{\mathcal O\text{ neither root orbit}}}
       m_{\mathcal O}s_{2b+1}(\mathcal O).
\]
Every coefficient displayed here is strictly positive. Equality
therefore permits only the short-root orbit and full-Cartan zero,
with $m_0=m_{\mathrm s}$. The highest weight must be $e_1$,
which gives the standard module. Its weights are
$0,\pm e_1,\ldots,\pm e_b$, once each, so equality is attained.
This reasoning also covers the half-integral weight lattice: those
weights have no zero coordinates and hence strictly positive surplus.

Now let $n+1=2b$ with $b\geq3$. The absolute root system is
$D_b$. Its even signed permutations give the same fraction
$j/(2b)$. If a zero coordinate is available, its sign can absorb the
parity restriction. If every coordinate is nonzero, flipping the
first and one other sign exchanges the two signs in the first
coordinate, so that case also has positive fraction $1/2$.
Thus
\[
 s_{2b}(\mathcal O)=(j-1)|\mathcal O|\geq0.
\]
The single root orbit has size $2b(b-1)$ and surplus
$2b(b-1)>b$. Its multiplicity $m_{\mathrm r}$ satisfies
$m_0\leq b m_{\mathrm r}$. Equation~\eqref{eq:panels-surplus}
therefore proves the bound. At equality the root orbit and
full-Cartan zero are absent, and every remaining weight has only
one nonzero coordinate. The highest weight is consequently
$a e_1$ with an integer $a\geq1$; half-integral weights cannot
have this pattern. If $a>1$, lowering by $e_1-e_2$ produces the
weight $(a-1)e_1+e_2$ by
Lemma~\ref{lem:panels-weight-multiplicities}. This has two
nonzero coordinates and positive surplus, a contradiction.
Hence $a=1$, giving the standard module with weights $\pm e_i$.

For completeness, in $D_2$ the two root orbits
$\{\pm(e_1-e_2)\}$ and $\{\pm(e_1+e_2)\}$ each have
size two, positive count one, and surplus two for $c=4$.
Each contains one absolute simple root, so each surplus strictly
exceeds its contribution to the zero-multiplicity bound. An explicit
calculation also handles all modules for the nonsimple complexification.
They are
\[
 \operatorname{Sym}^a\mathbb C^2\otimes
 \operatorname{Sym}^b\mathbb C^2,
 \qquad a,b\in\mathbb Z_{\geq0},\quad (a,b)\neq(0,0).
\]
Their split weights are
$(a+b-2i-2j)/2$ for $0\leq i\leq a$, $0\leq j\leq b$.
The dimension is $(a+1)(b+1)$, and the split-zero multiplicity is
zero if $a,b$ have different parity, and $\min(a,b)+1$ otherwise.
In the latter case assume $a\leq b$. The zero proportion is
$1/(b+1)\leq1/2$, with equality possible only when $a=b=1$:
if $b=0$ the representation is trivial, and if $b=1$ equal parity
forces $a=1$. Since the positive and negative weight counts agree,
\[
 \frac{\dim V}{h(V)}
 =\frac{2}{1-z(V)/\dim V}\leq4,
\]
with equality exactly for the four-dimensional standard module.

Finally, for $B_1$, or $SO(2,1)$, a covering-group irreducible has
highest split weight $a\in\frac12\mathbb Z_{>0}$ and the string
$a,a-1,\ldots,-a$. For integral $a$ its ratio is
$(2a+1)/a=2+1/a\leq3$, with equality only at $a=1$.
For strictly half-integral $a$ the zero space is absent and the
ratio is two. This proves the remaining case, including its spin
modules.
\end{proof}

\begin{lemma}[Unitary weight bound]\label{lem:panels-unitary-ratio}
For every nontrivial irreducible complex module $V$ of $SU(n,1)$,
with $k=n+1\geq3$, one has $\dim V\leq k h(V)$.
If $k>3$, equality holds precisely for $E,E^*$; if $k=3$, it
also holds for $\operatorname{Sym}^2E$ and
$\operatorname{Sym}^2E^*$.
\end{lemma}

\begin{proof}
Use the $A_{k-1}$ coordinates with sum zero and split functional
$\mu\mapsto\mu_1-\mu_2$. Suppose that the distinct coordinate
values of a nonzero weight have multiplicities
$r_1,\ldots,r_s$, where $s\geq2$ and $\sum r_i=k$.
Choosing the ordered first and second coordinates uniformly over a
permutation orbit gives
\[
 \frac{p(\mathcal O)}{|\mathcal O|}
 =\frac{k^2-\sum_i r_i^2}{2k(k-1)}\geq\frac1k.
\]
Indeed, merging parts increases their sum of squares, and among
two positive parts $r,k-r$ that sum is largest at $r=1$ or
$k-1$. Thus $\sum r_i^2\leq(k-1)^2+1$, with equality
exactly for the partition $(k-1,1)$. Equality in the orbit bound
therefore means that exactly one coordinate is exceptional.

The root orbit consists of the $k(k-1)$ vectors $e_i-e_j$.
There is one root of split weight two and $2(k-2)$ roots of split
weight one, so its surplus for $c=k$ is
\[
 k(2k-3)-k(k-1)=k(k-2)>k-1.
\]
If $m_{\mathrm r}$ is its multiplicity, then
$m_0\leq(k-1)m_{\mathrm r}$. All other orbit surpluses are
nonnegative. Equation~\eqref{eq:panels-surplus} proves the bound
and shows that equality forces $m_0=m_{\mathrm r}=0$ and only
one-exceptional-coordinate patterns.

A dominant highest weight with that pattern is
$a\omega_1$ or $a\omega_{k-1}$, where $a$ is a positive
integer. The corresponding modules are
$\operatorname{Sym}^aE$ and its dual. Their weights are obtained
from the monomial exponent tuples
$(a_1,\ldots,a_k)$, $a_i\geq0$, $\sum a_i=a$, by subtracting
$a/k$ from each coordinate. If $a\geq3$, the tuple
$(a-1,1,0,\ldots,0)$ has three distinct coordinate values, and
therefore positive surplus. If $a=2$ and $k>3$, the tuple
$(1,1,0,\ldots,0)$ has two groups of sizes $2,k-2$, again
giving positive surplus. These cases cannot attain equality.
For $a=1$, all weights have the required pattern and
$\dim E=k$, $h(E)=1$. For $a=2,k=3$, the only exponent
patterns are permutations of $(2,0,0)$ and $(1,1,0)$; both have
one exceptional coordinate and neither gives the full-Cartan zero
weight. In that case $\dim\operatorname{Sym}^2E=6$ and its
split weights are $2,1,0,0,-1,-2$, so $h=2$. Dualizing reverses
the split weights and preserves these conclusions.
\end{proof}

\begin{lemma}[Symplectic weight bound]\label{lem:panels-symplectic-ratio}
For every nontrivial irreducible complex module $V$ of $Sp(n,1)$,
put $k=n+1$. If $k\geq4$, then
$\dim V\leq k h(V)$, with equality exactly for $E$.
If $k=3$, then $\dim V\leq(7/2)h(V)$, with equality exactly
for $\bigwedge^3_0E$.
\end{lemma}

\begin{proof}
The absolute root system is $C_k$, its full weight lattice is
$\mathbb Z^k$, and the split functional is
$\mu\mapsto\mu_1+\mu_2$. The Weyl group consists of all signed
permutations. Suppose that a nonzero weight has $j$ nonzero
coordinates. The probability that just one of its first two
coordinates is nonzero is $2j(k-j)/(k(k-1))$. In that event the
split weight is positive with probability $1/2$. The probability
that both are nonzero is $j(j-1)/(k(k-1))$. Conditional on their
absolute values, at least one of the four sign choices has positive
sum; exactly one does if and only if those absolute values agree.
Consequently
\begin{equation}\label{eq:panels-symplectic-fraction}
 \frac{p(\mathcal O)}{|\mathcal O|}
 \geq\frac{j(k-j)+j(j-1)/4}{k(k-1)}.
\end{equation}
For $k\geq4$ the numerator is a strictly concave quadratic in
$j$. At the endpoints $j=1$ and $j=k$ its values are $k-1$
and $k(k-1)/4$, respectively. Thus the right side is at least
$1/k$. Equality can occur only for $j=1$, or for $k=j=4$;
in the latter case equality in the sign count further requires all
absolute values to agree.

Let $m_{\mathrm s}$ and $m_{\mathrm l}$ denote the
multiplicities of the short-root orbit
$\{\pm e_i\pm e_j:i\neq j\}$ and the long-root orbit
$\{\pm2e_i\}$. Lowering $2e_1$ by $e_1-e_2$ gives
$e_1+e_2$, so Lemma~\ref{lem:panels-weight-multiplicities}
gives $m_{\mathrm l}\leq m_{\mathrm s}$. Since there are
$k-1$ short absolute simple roots and one long one,
\[
 m_0\leq(k-1)m_{\mathrm s}+m_{\mathrm l}
       \leq k m_{\mathrm s}.
\]
The short-root orbit has size $2k(k-1)$. Its positive roots for
the split functional are $e_1+e_2$ and the $4(k-2)$ roots having
one index in $\{1,2\}$ and the other outside that set. Thus its
positive count is $4k-7$, and its surplus is
\[
 k(4k-7)-2k(k-1)=k(2k-5)>k.
\]
Together with Equation~\eqref{eq:panels-symplectic-fraction},
this proves the bound and shows that equality requires
$m_{\mathrm s}=m_{\mathrm l}=m_0=0$ and only zero-surplus
nonzero orbits. A highest weight of the first possible kind is
$a e_1$. If $a>1$, lowering by $e_1-e_2$ produces
$(a-1)e_1+e_2$, which has strictly positive surplus. The other
possible highest weight is $(a,a,a,a)$ when $k=4$. Lowering by
$e_3+e_4$ produces $(a,a,a-1,a-1)$. For $a=1$ this has
support two, and for $a>1$ it has unequal nonzero absolute values;
in either case its surplus is strictly positive. The only remaining
highest weight is $e_1$. The standard module has dimension $2k$
and two positive split weights, so it attains the bound.

For $C_3$, use $c=7/2$. The following table gives every possible
nonzero orbit pattern, its size, its positive count, and its surplus.
Letters denote positive integers, and distinct letters denote
distinct integers; permutations and all sign choices are included.
\[
\begin{array}{c|r|r|r}
 \text{pattern}&|\mathcal O|&p(\mathcal O)&
                    (7/2)p(\mathcal O)-|\mathcal O|\\ \hline
 (a,0,0)&6&2&1\\
 (a,a,0)&12&5&11/2\\
 (a,b,0)&24&12&18\\
 (a,a,a)&8&2&-1\\
 (a,a,b)&24&10&11\\
 (a,b,c)&48&24&36
\end{array}
\]
Here the first row has positive split weight exactly when its
nonzero coordinate occupies one of the first two positions and
has positive sign. In the second row, placing zero in the third
position contributes one positive sign choice, while placing it
in either of the first two positions contributes two, for a total
of five. Distinct nonzero absolute values in the third and sixth
rows prevent a zero sum of the first two coordinates, so exactly
half the orbit is positive. In the fourth row the first two signs
must both be positive, leaving two choices for the third sign.
In the fifth row, the exceptional absolute value occupies the
third position in eight orbit points, of which two are positive;
in the other sixteen points the first two absolute values differ,
so eight are positive. This gives all the entries without a
generic-position assumption on the integers.

Write $t_a$ for the multiplicity of the orbit $(a,a,a)$.
Lowering this weight by $e_1+e_2$ gives
$(a-1,a-1,a)$, whose orbit multiplicity is at least $t_a$.
For $a=1$ this companion orbit is $(1,0,0)$, with surplus one;
for $a>1$ it has pattern $(a,a-1,a-1)$, with surplus eleven.
The companion orbits are distinct as $a$ varies. None is the
short-root orbit $(1,1,0)$. Thus each negative contribution
$-t_a$ in Equation~\eqref{eq:panels-surplus} is paid for by a
distinct companion orbit; the payment leaves at least $10t_a$
when $a>1$.

Independently, the earlier lowering comparison of long and short
roots gives $m_0\leq3m_{\mathrm s}$. The short-root surplus
is $11/2>3$, so its contribution pays for full-Cartan zero with
strict excess unless $m_{\mathrm s}=m_0=0$. Every other
unassigned surplus is positive. The bound follows. At equality,
all $t_a$ with $a>1$ vanish, and the only remaining orbits can be
$(1,1,1)$ and $(1,0,0)$ with equal multiplicity. The module is
nontrivial, so both occur. Since $(1,1,1)-(1,0,0)=e_2+e_3$
is a positive root, its highest weight must be $(1,1,1)$.
Its highest-weight multiplicity is one, so both orbit
multiplicities are one.

This module is $\bigwedge^3_0E$. Indeed, the standard weights
are $\pm e_1,\pm e_2,\pm e_3$. In $\bigwedge^3E$,
choosing one signed weight of each index gives all eight
$(\pm1,\pm1,\pm1)$ weights once. Choosing an opposite pair
and one other signed weight gives each $\pm e_i$ twice,
because either of the other two indices can supply the opposite
pair. Hence $\bigwedge^3E$ has these eight weights and twelve
sparse weight vectors. The invariant symplectic bivector gives
an injective equivariant map $E\to\bigwedge^3E$,
$v\mapsto\Omega\wedge v$; on a symplectic basis vector its
image is a sum of two distinct nonzero wedge-basis vectors.
Removing this standard summand leaves the eight sign weights
and the six sparse weights, each once. To verify irreducibility
without assuming it, the possible dominant highest weights in
this complement are $(1,1,1)$ and $(1,0,0)$. The former occurs
once and, by lowering $e_1+e_2$, its irreducible constituent
already contains the sparse orbit. It therefore contains all
fourteen available weight vectors. The complement is irreducible,
has dimension fourteen and positive count four, as required.
\end{proof}

The following argument uses convolution of orbital measures to force
normalized-character decay, a method developed in
\cite{Ragozin1972,HareWilsonYee2000}. We give the form needed for the
uniform gap in the present representation count.

\begin{lemma}[Decay of the zero-weight proportion]
\label{lem:panels-zero-weight-decay}
Let $K_c$ be a compact connected semisimple Lie group, and let
$A\in\mathfrak k_c$ have nonzero projection to every simple ideal.
Suppose $\exp(TA)=1$ for some $T>0$. For its irreducible
representations $\pi$, of dimension $d_\pi$, one has
\[
 \frac{\dim\ker d\pi(A)}{d_\pi}\longrightarrow0
\]
as $\pi$ leaves every finite subset of the irreducible dual.
\end{lemma}

\begin{proof}
First fix $x\in K_c$ such that $\Ad x$ is not the identity on
any simple ideal. Put
\[
 U=(1-\Ad x)\mathfrak k_c.
\]
The span of $\Ad(k)U$ over $k\in K_c$ is an
$\Ad(K_c)$-invariant ideal. Its projection to every simple
ideal is nonzero, and hence it is all of $\mathfrak k_c$.
Choose $h_1,\ldots,h_N\in K_c$ so that
$\sum_j\Ad(h_j)U=\mathfrak k_c$. Consider the analytic map
\[
 F:K_c^N\longrightarrow K_c,
 \qquad F(k_1,\ldots,k_N)=\prod_{j=1}^N k_jxk_j^{-1}.
\]
Write $p_0=1$ and
$p_j=\prod_{\ell\leq j}k_\ell xk_\ell^{-1}$.
For variations $k_j(t)=k_j\exp(tX_j)$, right trivialization
of the differential gives
\[
 dF(X_1,\ldots,X_N)F^{-1}
 =\sum_{j=1}^N
   \Ad(p_{j-1}k_j)(1-\Ad x)X_j.
\]
Choose $k_j=p_{j-1}^{-1}h_j$ recursively. At that point the
displayed differential is onto.

Equip the groups with invariant analytic metrics. In invariant
frames the matrix $B$ of $dF$ is analytic, and
$\det(BB^{\mathsf T})$ is an analytic function on the connected
manifold $K_c^N$, nonzero at the constructed point. Its zero
set has Haar measure zero. On the remaining set, submersion
coordinates make $F$ a coordinate projection. A countable
cover by such coordinates and Fubini's theorem show that the
inverse image of any Haar-null subset of $K_c$ is null there.
The critical set is also null. Therefore the pushforward of
product Haar measure under $F$ is absolutely continuous.
Equivalently, if $\nu_x$ is the conjugacy-orbit probability
measure of $x$, then
\[
 \nu_x^{*N}=f_x\,dk,\qquad f_x\in L^1(K_c).
\]

For an irreducible unitary $\pi$, Schur's lemma gives
\[
 \int\pi(g)\,d\nu_x(g)
 =a_\pi I,\qquad a_\pi=\frac{\chi_\pi(x)}{d_\pi}.
\]
Taking normalized traces after $N$ convolutions yields
\[
 a_\pi^N=\int_{K_c}f_x(g)
                    \frac{\chi_\pi(g)}{d_\pi}\,dk.
\]
Approximate $f_x$ in $L^1$ by a function $f_{x,M}\in L^2$.
Orthogonality of irreducible characters and Bessel's inequality
give $\int f_{x,M}\chi_\pi\,dk\to0$ as $\pi$ leaves
finite subsets; the conjugate characters form the same
orthonormal family. Moreover,
\[
 \left|\int(f_x-f_{x,M})\frac{\chi_\pi}{d_\pi}\,dk\right|
 \leq\|f_x-f_{x,M}\|_1,
\]
because $|\chi_\pi|\leq d_\pi$. First letting $\pi$ leave
finite subsets and then letting $M$ tend to infinity proves
$a_\pi^N\to0$, and therefore $a_\pi\to0$.

For each simple ideal choose a nonzero eigenvalue $i\beta$ of
$\operatorname{ad}A$ on that ideal. Such an eigenvalue exists
because the ideal has zero center and the projection of $A$ is
nonzero. If $\exp(tA)$ acts centrally on that ideal, then
$t\beta\in2\pi\mathbb Z$. The union of these exceptional
parameter sets, for $0\leq t\leq T$, has measure zero.
For every other $t$, the preceding normalized-character decay
applies. Averaging a representation over the periodic subgroup
projects onto $\ker d\pi(A)$, so
\[
 \frac{\dim\ker d\pi(A)}{d_\pi}
 =\frac1T\int_0^T
      \frac{\chi_\pi(\exp(tA))}{d_\pi}\,dt.
\]
The integrands are bounded in absolute value by one.
Dominated convergence proves the assertion along every sequence
leaving finite subsets, hence on the entire irreducible dual.
The argument uses all simple ideals at once and therefore also
applies when $K_c$ has several simple factors.
\end{proof}

\begin{lemma}[Invariant forms on maximizing modules]
\label{lem:panels-maximizer-forms}
Every maximizing module in Lemmas
\ref{lem:panels-orthogonal-ratio}--\ref{lem:panels-symplectic-ratio}
has the weight, form, and parabolic-stage properties in
Lemma~\ref{lem:panels-optimal-modules}.
\end{lemma}

\begin{proof}
We use a positive Hilbert norm invariant under a compact real form
of the complexified group. In each standard realization, the
split generator is Hermitian for this norm. Write an invariant
Hermitian form as $q(v,w)=\langle Qv,w\rangle$, using
consistent conventions for the two Hermitian products. Its
invariance under $\exp(tH_R)$ implies
\[
 d\pi(H_R)^*Q+Qd\pi(H_R)=0.
\]
When the split generator is Hermitian, this says that $Q$ takes
$V_t$ to $V_{-t}$. Nondegeneracy therefore pairs opposite
weight spaces perfectly and all other pairs trivially. Each pair
$V_t\oplus V_{-t}$ for $t>0$ has signature
$(\dim V_t,\dim V_t)$. It follows that if the negative index
of $q$ is $h(V)$, its restriction to $V_0$ is positive
definite. We verify that index and the normalization $Q^2=I$
in all the cases.

For the orthogonal standard module of dimension $N=n+1$,
take the Hermitian extension of the real quadratic form of
signature $(N-1,1)$ and the Euclidean compact-form norm.
Its symmetry is $Q=\operatorname{diag}(1,\ldots,1,-1)$.
A split generator interchanges one positive basis vector and
the negative basis vector, with zero action on their orthogonal
complement. Its eigenvalues are $1,-1$ and $N-2$ zeros.
Thus $h=1$ is the negative index and the zero space is
positive and nonzero. This includes the orthogonal low-rank
names used above.

For the unitary standard module of dimension $k$, use the
Hermitian form of signature $(k-1,1)$ and its usual positive
compact-form norm. The same two-dimensional boost gives
split weights $1,-1$, and zero with multiplicity $k-2$.
Again the negative index is $h=1$. On the dual, use the
inverse Hermitian form and the dual Hilbert norm; its matrix
is the transpose inverse of $Q$, still a self-adjoint
involution, with the same signature and opposite split weights.

For $SU(2,1)$, the induced form on
$\operatorname{Sym}^2E$ has positive index
$\dim\operatorname{Sym}^2\mathbb C^2+1=4$ and negative
index two: the negative vectors are the products of one vector
from the positive two-plane with the negative basis vector.
The induced symmetry is the restriction of $Q\otimes Q$
to the symmetric tensors, and hence is a self-adjoint
involution for the induced compact-form norm. The induced
split generator is Hermitian. Its split spectrum is
$2,1,0,0,-1,-2$, so $h=2$ equals the negative index and
the two-dimensional zero space is positive. The dual argument
gives the identical conclusions for $\operatorname{Sym}^2E^*$.

For the standard $Sp(n,1)$ module, of complex dimension
$2k=2n+2$, group the orthonormal basis into symplectic pairs
$e_i,f_i$. Give each of the first $k-1$ pairs sign $+1$ and
the last pair sign $-1$. The preserved Hermitian form has
signature $(2k-2,2)$, its symmetry is diagonal with these
signs, and a split generator couples one positive pair to the
negative pair. It has weights $1,-1$, each twice, and zero
with multiplicity $2k-4$. Thus its negative index is $h=2$
and its zero space is positive. These are integral weights.

It remains to examine the primitive cube for $Sp(2,1)$.
Let $E=\mathbb C^6$ with compact-orthonormal basis
$e_i,f_i$, $1\leq i\leq3$, and let
\[
 q(e_i,e_j)=q(f_i,f_j)=s_i\delta_{ij},
 \qquad (s_1,s_2,s_3)=(1,1,-1),
 \qquad \omega(e_i,f_j)=\delta_{ij}.
\]
All other pairings between distinct displayed basis vectors
vanish, with $\omega$ extended skew-symmetrically. Write
$C e_i=s_i e_i$ and $C f_i=s_i f_i$. The Hermitian form
$q_3$ on $\bigwedge^3E$ has signature
\[
 \left(\binom43+4\binom22,\ 2\binom42\right)=(8,12).
\]
The first number counts wedges with zero or two negative
basis vectors, and the second counts wedges with one.
Put $\Omega=\sum_i e_i\wedge f_i$ and
$L(v)=\Omega\wedge v$. For each basis vector, $L(v)$
contains the two complete pairs of the other indices.
Each term has signed squared norm equal to that of $v$,
since the squared sign of a complete pair is one. The wedge
basis supports for distinct unpaired basis vectors are disjoint.
Consequently
\[
 q_3(Lv,Lw)=2q(v,w).
\]
In particular $L$ is injective and its image has signature
$(4,2)$.

Let $c:\bigwedge^3E\to E$ be symplectic contraction,
normalized as the adjoint of $L$ for the positive compact-form
norm. Explicitly it contracts each complete $e_i,f_i$ pair
with its wedge sign; equivalently,
\[
 c(u\wedge v\wedge w)
 =\omega(u,v)w-\omega(u,w)v+\omega(v,w)u.
\]
This formula shows equivariance under the complex symplectic
group, and a basis calculation gives $cL=2I$. Since
$(\bigwedge^2C)\Omega=\Omega$, one has
\[
 (\bigwedge^3C)L=LC.
\]
It follows both that $L(E)$ reduces the self-adjoint involution
$\bigwedge^3C$ and that
$q_3(Lv,w)=q(v,cw)$. Therefore
\[
 P=\bigwedge^3_0E=\ker c
   =L(E)^{\perp_{q_3}}
\]
is also the compact-form orthogonal complement of $L(E)$.
Its signature for $q_3$ is $(4,10)$. On $P$ choose instead
$q_P=-q_3|_P$, whose signature is $(10,4)$ and whose
symmetry is $-(\bigwedge^3C)|_P$, a self-adjoint involution.
The compact-form orthogonal decomposition is invariant under
the compact real form, hence under its complexified Lie algebra;
the split generator remains Hermitian on $P$.

The weight calculation in
Lemma~\ref{lem:panels-symplectic-ratio} gives the primitive
split spectrum
\[
\begin{array}{c|rrrrr}
 t&2&1&0&-1&-2\\ \hline
 \dim P_t&2&2&6&2&2.
\end{array}
\]
Thus $h(P)=4$ is exactly the negative index of $q_P$;
the general opposite-pair argument already proves positivity
on the six-dimensional $P_0$. One can also see its sign
directly. Choose the split generator
\[
 H_Re_2=e_3,\quad H_Re_3=e_2,\qquad
 H_Rf_2=-f_3,\quad H_Rf_3=-f_2,\qquad
 H_Re_1=H_Rf_1=0.
\]
It is Hermitian for the compact norm, preserves $\omega$
infinitesimally, and anticommutes with $C$. Define
\[
 a_1=(e_2+e_3)/\sqrt2,\quad
 a_2=(f_2-f_3)/\sqrt2,\qquad
 b_1=(e_2-e_3)/\sqrt2,\quad
 b_2=(f_2+f_3)/\sqrt2.
\]
The $a_i$ have split weight one, the $b_i$ have split
weight minus one, and $q(a_i,b_j)=\delta_{ij}$ while
each of their two spans is isotropic. For $z=e_1,f_1$ put
$w_{ij}(z)=a_i\wedge z\wedge b_j$. These eight vectors
span the zero space in the full cube and satisfy
\[
 q_3(w_{ij}(z),w_{kl}(z'))
       =-\delta_{il}\delta_{jk}\delta_{zz'}.
\]
For each $z$, the positive direction is $w_{12}-w_{21}$;
the three negative directions are $w_{11},w_{22}$ and
$w_{12}+w_{21}$. Moreover,
\[
 \Omega=e_1\wedge f_1+a_1\wedge b_2-a_2\wedge b_1,
 \qquad L(z)=-(w_{12}(z)-w_{21}(z)).
\]
Removing $L(E)_0$ therefore removes exactly the two positive
directions, leaving six negative directions for $q_3$, as
asserted. Negating the form makes all of them positive.

In every case all listed weights are integral and the zero
spaces just computed are nonempty. The statements about
annihilators now follow from the perfect opposite-weight
pairings and the positive zero space. To verify the parabolic
assertion, use a positive multiple of $H_R$ that is an
integral cocharacter of the simply connected complexified
group. A root space whose root is nonnegative on $H_R$
maps a weight space $V_t$ into $V_{t+\alpha(H_R)}$ and
therefore preserves both $V_{>0}$ and $V_{\geq0}$.
The corresponding algebraic parabolic, generated by this
Cartan and these root groups, preserves both spaces. This is
the complex parabolic attached to the oriented rank-one
boundary point. Its translates are algebraic subspaces in
the appropriate Grassmannians, hence give the asserted
parabolic stages independently of the choice of opposite
point. Their middle quotient identifies with $V_0$ and
inherits its positive form.
\end{proof}

\begin{proof}[Proof of Lemma~\ref{lem:panels-optimal-modules}]
The bounds and their complete equality lists follow from
Lemmas~\ref{lem:panels-orthogonal-ratio},
\ref{lem:panels-unitary-ratio}, and
\ref{lem:panels-symplectic-ratio}; the form and integrality
claims follow from Lemma~\ref{lem:panels-maximizer-forms}.
The displayed spectra give $d_R=2$ in the two stated
exceptional cases and $d_R=1$ in the others.

It remains to obtain one uniform gap for all other
irreducibles. Let $K_c$ be the simply connected compact
real form of the complexified Lie algebra and set $A=iH_R$
in that compact real form. In the coordinate realizations
used above, $A$ has a period after passing to the full
covering-group weight lattice. Its projection to every
simple ideal is nonzero. This is immediate in the simple
cases; for $SO(3,1)$ the two $\mathfrak{sl}_2$ components
of the displayed split functional are both nonzero.
Lemma~\ref{lem:panels-zero-weight-decay} consequently gives
$z(V)/\dim V\to0$ as the highest weight leaves finite
subsets. The split functional and its negative are Weyl
conjugate in each of these realizations, so
\[
 h(V)=\frac{\dim V-z(V)}2,
 \qquad \frac{\dim V}{h(V)}
       =\frac2{1-z(V)/\dim V}\longrightarrow2.
\]
Every tabulated $c_R$ is greater than two. Choose
$\tau$ with $2<\tau<c_R$. Only finitely many
irreducibles have ratio at least $\tau$. After the
maximizing modules are removed, the ratios of this finite
set are all strictly smaller than $c_R$ by the equality
classification. The maximum of those remaining ratios
and $\tau$ is therefore a number $b_R<c_R$.
Taking $\varepsilon_R=c_R-b_R$ proves the uniform
gap. Finally, reversing $H_R$ interchanges positive and
negative spaces, preserves all ratios, and preserves the
largest absolute weight. Thus $c_R$ and $d_R$ are
unchanged at the opposite orientation.
\end{proof}

The representation-theoretic verification completes the real-panel upper
estimate.  Section~\ref{sec:heights} combines it with the other local
branches and uses Haar volume to calibrate these one-sided bounds.
